% ========================================
% COMANDOS PARA FORMATAÇÃO DA TESE UNESP
% ========================================

% ----- DEFINIÇÃO DOS AMBIENTES PRÉ-TEXTUAIS -----
\newenvironment{dedicatoria}{%
    \cleardoublepage
    \thispagestyle{empty}
    \vspace*{\fill}
    \begin{flushright}
    \itshape
}{%
    \end{flushright}
    \vspace*{\fill}
    \cleardoublepage
}

\newenvironment{agradecimentos}{%
    \cleardoublepage
    \thispagestyle{empty}
    \chapter*{AGRADECIMENTOS}
}{%
    \cleardoublepage
}

\newenvironment{epigrafe}{%
    \cleardoublepage
    \thispagestyle{empty}
    \vspace*{\fill}
    \begin{flushright}
    \itshape
}{%
    \end{flushright}
    \vspace*{\fill}
    \cleardoublepage
}

\newenvironment{resumo}{%
    \cleardoublepage
    \thispagestyle{empty}
    \chapter*{RESUMO}
    \vspace{0.5cm}
}{%
    \cleardoublepage
}

\newenvironment{abstract}{%
    \cleardoublepage
    \thispagestyle{empty}
    \chapter*{ABSTRACT}
    \vspace{0.5cm}
}{%
    \cleardoublepage
}

% ----- COMANDO PARA FIGURAS -----
\newcommand{\figura}[4]{%
    \begin{figure}[H]
        \centering
        \includegraphics[width=#1]{#2}
        \caption{#3}
        \label{fig:#4}
    \end{figure}
}

% ----- COMANDO PARA TABELAS -----
\newcommand{\tabela}[4]{%
    \begin{table}[H]
        \centering
        \caption{#1}
        \label{tab:#2}
        #3
        \fonte{#4}
    \end{table}
}

\newcommand{\fonte}[1]{%
    \vspace{0.2cm}
    \par\small\textbf{Fonte:} #1
}

\newcommand{\nota}[1]{%
    \vspace{0.2cm}
    \par\small\textbf{Nota:} #1
}